% !TeX program = xelatex
% 40分钟新课默认模板。整目录复制；时长写在注释及教师教案中。
\documentclass[aspectratio=169,10pt]{beamer}
\usepackage[UTF8,fontset=fandol]{ctex}
\IfFontExistsTF{Microsoft YaHei}{\setCJKsansfont{Microsoft YaHei}}{}
\usetheme{GaoZhongYuan}
\title{课题}\author{授课教师}\date{}
\begin{document}\campusmaketitle
% 0—3分钟：情境引入。全课探究讲解共27分钟，包含学生思考和即时练习。
\begin{frame}[t]{问题引入}
  写一个能暴露本课认知需要的问题，给学生思考时间。
\end{frame}
% 3—10分钟：概念探究与条件。
\begin{frame}[t]{概念与条件}
  呈现定义、适用条件，引导学生解释关键条件。
\end{frame}
% 10—16分钟：示范与即时检查。每题、每个小问独占一页。
\begin{campuspractice}{例题}
  在这里写一道用于解释概念的题目。
\end{campuspractice}
\begin{campuspractice}{即时检查}
  在这里写一道同结构但能检验理解的题目。
\end{campuspractice}
% 16—21分钟：性质与依据。
\begin{frame}[t]{性质与依据}
  由定义或已知事实推导性质，写清推理依据。
\end{frame}
% 21—27分钟：应用与再次检查。
\begin{campuspractice}{应用例题}
  在这里写一道关键方法例题。
\end{campuspractice}
\begin{campuspractice}{即时检查}
  在这里写一道可暴露易错点的检查题。
\end{campuspractice}
% 27—33分钟：集中独立练习。4个短小问是本课示例，可按难度调整。
\begin{campuspractice}{当堂练习1}
  检查概念或基本表示。
\end{campuspractice}
\begin{campuspractice}{当堂练习2}
  检查基本方法。
\end{campuspractice}
\begin{campuspractice}{当堂练习3}
  检查适用条件。
\end{campuspractice}
\begin{campuspractice}{当堂练习4}
  检查容易遗漏的情形。
\end{campuspractice}
% 33—37分钟：依据巡视、板演反馈。参考答案仅教师版可见，勿覆盖演算区。
\ifdefined\TeacherEdition
\begin{frame}[t]{当堂练习：反馈}
  在这里写参考答案、错因与纠正依据。
\end{frame}
\fi
% 37—40分钟：学生先回顾，教师再补齐条件与方法。
\begin{campussummary}
\begin{itemize}
\item 本课核心概念是什么？
\item 使用它需要哪些条件？
\item 关键方法的依据是什么？
\item 怎样检查本课最容易出现的错误？
\end{itemize}
\end{campussummary}
% 全量题库、拓展材料另存；不顺序接在主课后默认全讲。
\campusclosing[谢谢]\end{document}
